\honest{A Priori}{Eliezer Yudkowsky}{2007}

Traditional Rationality treats its rules as social rules; to a ``Bayesian,'' the brain is an engine that stops working when they are broken.

Why believe the simpler of two hypotheses that fit the facts equally well? Not because Occam's razor has worked before: the hypothesis that it stops working after today, the day I post this, fits the same record. Not because a simplicity prior is ``reasonable'': who says which prior is reasonable? It seems the razor can be justified only by the razor. \nb{This is Hume's circularity argument about induction, and the hypothesis that the razor stops working after today has the form of Goodman's ``grue''.}

A philosopher may shrug, call a truce with colleagues, and fly the white flag of ``a priori truth.'' I name no one. \nb{The Stanford Encyclopedia's entry on simplicity shares part of the complaint: with a priori defences of simplicity, it can be hard to tell ``an a priori defense'' from ``no defense''. It also reports that many philosophers had turned away from such defences by the late twentieth century.} If there were an a priori truth factory, why could a thirsty hunter-gatherer not use it to find water? \nb{The standard examples of a priori knowledge are arithmetic, logic and conceptual truths. Nothing in them concerns where to find water.}

Hume defined these truths as ``discoverable by the mere operation of thought, without dependence on what is anywhere existent in the universe.'' But thoughts exist; they are brain events. So when pure thought tells you that 1 + 1 = 2, ``you are, in effect, observing your own brain as evidence,'' and you would learn exactly the same by watching another brain do the sum. ``There is nothing you can know `a priori', which you could not know with equal validity'' by observation. \nb{Hume's next sentence shows that the independence he meant was from the objects, circles and triangles, not from the thinker. The step from how a belief is produced to what justifies it needs the view called reliabilism, which the post does not name. Reliabilists have also used it to defend the a priori, and on one view of testimony a result learned from another mind keeps its a priori justification. Whether anything is known a priori was, and is, debated among philosophers.}

The brain's bias toward simplicity is not arbitrary; regularized algorithms work better. Occam's razor works because the world is simple; why it is simple, ``This I do not know.'' A mind without modus ponens cannot be argued into it, any more than a rock. \nb{That is the point of Lewis Carroll's ``What the Tortoise Said to Achilles'' (1895), which Yudkowsky discussed in a later post.} Brains were built by selection, not argued into existence. To be satisfied with this, ``one must see rationality in terms of engines, rather than arguments.'' \nb{This is the externalist reply to Hume's circularity: a rule that is in fact reliable can yield justified beliefs without first being justified.}
